\%\textbackslash documentclass{[}reprint,amsmath,amssymb,aps,twoside{]}\{revtex4-2\}
\% for editor use only

\% For initial submission use these lines
\textbackslash documentclass{[}amsmath,amssymb,aps, endfloats,
linenumbers{]}\{revtex4-2\}

\% DO NOT CHANGE THINGS BELOW HERE \textbackslash usepackage\{graphicx\}
\textbackslash usepackage\{amsmath,amssymb,amsfonts\}
\textbackslash usepackage\{dcolumn\} \textbackslash usepackage\{bm\}
\textbackslash usepackage\{siunitx\}
\%\textbackslash usepackage\{tikz,pgfplots\}
\textbackslash sisetup\{separate-uncertainty=true,
multi-part-units=single\}
\textbackslash usepackage{[}colorlinks,allcolors=blue{]}\{hyperref\}
\textbackslash usepackage\{cleveref\}
\textbackslash crefname\{equation\}\{\}\{\}
\textbackslash crefname\{figure\}\{Fig.\}\{Figs.\}
\textbackslash crefname\{table\}\{Table\}\{Tables\}
\textbackslash usepackage\{svg\} \textbackslash svgpath\{\{./figures\}\}

\% AUTHORS EDIT THIS PART \textbackslash hypersetup\{\% pdftitle=\{The
Effect of Popper Design on Maximum Jump Height\}, \% full title here
pdfauthor=\{Kaitlin Coulanges, Saanvi Dakwale, Samantha Hein, Reese
Wallace, and Sashank Yellapragada\}, \% full author list here \}

\textbackslash usepackage\{fancyhdr\} \textbackslash pagestyle\{fancy\}
\textbackslash fancyhf\{\} \textbackslash fancyhead{[}RE,RO{]}\{J
S\textbackslash\&E \textbackslash textbf\{2\}, xx-\/-xx (2026)\}
\textbackslash fancyhead{[}LO{]}\{Yellapragada \textbackslash emph\{et
al\}\} \textbackslash fancyhead{[}LE{]}\{The Effect of Popper Design on
Maximum Jump Height\}
\textbackslash fancyfoot{[}C{]}\{\textbackslash thepage\}
\textbackslash fancypagestyle\{mytitlepage\}\{
\textbackslash fancyhf\{\} \textbackslash fancyhead{[}C{]}\{Journal of
Science \textbackslash\& Engineering \textbackslash textbf\{2\},
xx-\/-xx (2026)\} \% FOR EDITOR ONLY
\textbackslash fancyfoot{[}C{]}\{\textbackslash thepage\} \}
\textbackslash begin\{document\}

\% AUTHORS EDIT THIS PART \textbackslash title\{The Effect of Popper
Design on Maximum Jump Height\}

\textbackslash author\{Kaitlin Coulanges\} \textbackslash author\{Saanvi
Dakwale\} \textbackslash author\{Samantha Hein\}
\textbackslash author\{Reese Wallace\} \textbackslash author\{Sashank
Yellapragada\} \textbackslash email\{Contact author:
427syellapragada@frhsd.com\} \% Replace with the lead
contact\textquotesingle s email

\textbackslash affiliation\{Science \textbackslash\& Engineering Magnet
Program, \textbackslash href\{https://manalapan.frhsd.com/\}\{Manalapan
High School\}, Englishtown, NJ 07726 USA\}
\textbackslash date\{\textbackslash today\}

\textbackslash begin\{abstract\} This experiment investigates whether
the facial expression of a popper affects its maximum jump height using
principles of energy conservation. Five
``happy\textquotesingle\textquotesingle{} poppers and five
``sad\textquotesingle\textquotesingle{} poppers were compressed to
maximum depth and released, converting elastic potential energy into
kinetic energy, and then gravitational potential energy. At peak height,
mean GPE for happy poppers was \textbackslash num\{96.2
\textbackslash pm 15.6\}
\textbackslash unit\{\textbackslash milli\textbackslash joule\} and for
sad poppers was \textbackslash num\{83.1 \textbackslash pm 9.2\}
\textbackslash unit\{\textbackslash milli\textbackslash joule\}. A
two-sample, two-tailed \$t\$-test revealed no statistically significant
difference in peak heights (\$p = 0.24 \textgreater{}
\textbackslash alpha = 0.05\$), or in masses (\$p = 0.12 \textgreater{}
\textbackslash alpha = 0.05\$). These results support our hypothesis
that facial expression does not have any meaningful effect on energy
transfer efficiency, and is dictated mainly by elastic properties.
\textbackslash end\{abstract\}

\textbackslash keywords\{kinematics, energy conservation, elastic
potential energy, poppers\}

\textbackslash maketitle\textbackslash thispagestyle\{mytitlepage\}
\textbackslash section\{Introduction\} When a popper is compressed, work
is done to deform the material, storing elastic potential energy. This
energy can be expressed as: \textbackslash begin\{equation\} U\_s =
\textbackslash frac\{1\}\{2\}kx\^{}2 \textbackslash label\{eq:epe\}
\textbackslash end\{equation\} where \$k\$ is the spring constant and
\$x\$ is the compression distance \textbackslash cite\{tipler\}. Upon
release, this stored elastic potential energy is rapidly converted into
kinetic energy (\$K\$), given by: \textbackslash begin\{equation\} K =
\textbackslash frac\{1\}\{2\}mv\^{}2 \textbackslash label\{eq:ke\}
\textbackslash end\{equation\} where \$m\$ is the mass and \$v\$ is the
velocity \textbackslash cite\{tipler\}. As the popper rises, its kinetic
energy is transformed into gravitational potential energy (GPE) and
represented by the symbol \$U\_g\$. At the maximum height (\$h\$), the
popper's velocity is approximately zero, meaning the energy is
momentarily stored as \$U\_g\$, described as:
\textbackslash begin\{equation\} U\_g = mgh
\textbackslash label\{eq:gpe\} \textbackslash end\{equation\} where
\$g\$ is the acceleration due to gravity \textbackslash cite\{tipler\}.

By measuring the maximum height reached, one can estimate the efficiency
of this energy transfer. Theoretically, a greater maximum height
indicates a greater conversion of elastic potential energy into upward
motion. This experiment investigates whether the design of a popper,
specifically a ``happy\textquotesingle\textquotesingle{} versus a
``sad\textquotesingle\textquotesingle{} facial expression, influences
the maximum height reached. We hypothesize that the facial expression
will not significantly affect the height, as energy transfer primarily
relies on the physical mass and the elastic properties of the popper.

\textbackslash section\{Methods and materials\} This experiment was
conducted using commercially available rubber poppers (Liberty Imports;
Shantou, CN) with multiple facial expressions. Two types were chosen:
``happy,\textquotesingle\textquotesingle{} identified by an
upward-curved mouth, and ``sad,\textquotesingle\textquotesingle{}
denoted by a downward-curved mouth and tear-like markings beneath each
eye. A sample size of \$n=5\$ was used for each facial expression group.
The mass of each popper was measured on a digital scale (Smart Weight;
Zhongshan, China). All trials were conducted indoors on a flat, level
surface within a classroom to maintain consistent environmental
conditions. For measuring jump height, two meter sticks (Westcott Rule
Company) were taped vertically to the wall. For each trial, to reduce
experimental error, the popper was compressed to its maximum depth on
the floor. One student launched all the poppers to avoid experimental
variability. To ensure accuracy, the maximum height of each jump was
documented on an iPhone 16 (Apple Inc; Cupertino, CA), recording at
\textbackslash qty\{240\}\{fps\} with \textbackslash num\{1080\}p and a
standard wide-angle lens, placed at a fixed height above the ground.
Specific landmarks of surrounding signs and posters were used to
calibrate height of the popper after launch, from which jump heights for
each popper were determined and recorded. To reduce bias, the order of
the trials was randomized rather than testing all poppers of one facial
expression. Trials in which the popper failed to launch properly were
discarded. Outliers, defined as trials where the height was two or more
standard deviations from the mean, were also excluded from the final
analysis. Mass and jump height were analyzed using a Google Sheets
\textbackslash cite\{googlesheets\} to calculate the mean and standard
deviation for each group. \textbackslash texttt\{R\}
\textbackslash cite\{R\}, with library \textbackslash texttt\{ggplot2\}
\textbackslash cite\{ggplot2\}, was used to create a bar chart comparing
the mean maximum jump height for happy and sad poppers. Two-sample,
two-tailed \$t\$-tests were performed to compare both the masses and the
maximum jump heights of the happy and sad poppers. A two-sample,
two-tailed \$t\$-test was also performed. The following hypotheses were
tested:

\textbackslash begin\{itemize\} \textbackslash item
\textbackslash textbf\{Masses:\} \$H\_\{0,A\}:
\textbackslash mu\_\{\textbackslash text\{happy\}\} =
\textbackslash mu\_\{\textbackslash text\{sad\}\}\$ vs. \$H\_\{1,A\}:
\textbackslash mu\_\{\textbackslash text\{happy\}\} \textbackslash neq
\textbackslash mu\_\{\textbackslash text\{sad\}\}\$ \textbackslash item
\textbackslash textbf\{Heights:\} \$H\_\{0,B\}:
\textbackslash mu\_\{\textbackslash text\{happy\}\} =
\textbackslash mu\_\{\textbackslash text\{sad\}\}\$ vs. \$H\_\{1,B\}:
\textbackslash mu\_\{\textbackslash text\{happy\}\} \textbackslash neq
\textbackslash mu\_\{\textbackslash text\{sad\}\}\$ \textbackslash item
\textbackslash textbf\{GPE:\} \$H\_\{0,C\}:
\textbackslash mu\_\{\textbackslash text\{happy\}\} =
\textbackslash mu\_\{\textbackslash text\{sad\}\}\$ vs. \$H\_\{1,C\}:
\textbackslash mu\_\{\textbackslash text\{happy\}\} \textbackslash neq
\textbackslash mu\_\{\textbackslash text\{sad\}\}\$
\textbackslash end\{itemize\}

The experimental setup is illustrated in
\textbackslash cref\{fig:setup\}.

\textbackslash begin\{figure\}{[}ht{]} \textbackslash centering
\textbackslash includegraphics{[}width=0.5\textbackslash textwidth{]}\{figures/setupphoto.png\}
\textbackslash caption\{Experimental setup showing the vertical
calibration scale and popper placement.\}
\textbackslash label\{fig:setup\} \textbackslash end\{figure\}
\textbackslash section\{Results\} The measured mass, maximum jump
height, and calculated gravitational potential energy (GPE) for the
``happy\textquotesingle\textquotesingle{} and
``sad\textquotesingle\textquotesingle{} popper groups are summarized in
\textbackslash cref\{tab:happy\} and \textbackslash cref\{tab:sad\},
respectively.

\textbackslash begin\{table\}{[}ht{]} \textbackslash caption\{Mass,
height, and potential energy for Happy Poppers.\}
\textbackslash label\{tab:happy\} \textbackslash begin\{ruledtabular\}
\textbackslash begin\{tabular\}\{cccc\} Trial \& Mass
(\textbackslash unit\{\textbackslash kilo\textbackslash gram\}) \&
Height (\textbackslash unit\{\textbackslash meter\}) \& GPE
(\textbackslash unit\{\textbackslash milli\textbackslash joule\})
\textbackslash\textbackslash{} \textbackslash colrule 1 \& 0.0058 \&
2.12 \& 120.5 \textbackslash\textbackslash{} 2 \& 0.0057 \& 1.82 \&
101.7 \textbackslash\textbackslash{} 3 \& 0.0056 \& 1.47 \& 80.8
\textbackslash\textbackslash{} 4 \& 0.0057 \& 1.67 \& 93.4
\textbackslash\textbackslash{} 5 \& 0.0056 \& 1.54 \& 84.6
\textbackslash\textbackslash{} \textbackslash colrule Mean \&
\textbackslash num\{0.0057\} \& \textbackslash num\{1.72\} \&
\textbackslash num\{96.2\} \textbackslash\textbackslash{} Std. Dev. \&
\textbackslash num\{0.00008\} \& \textbackslash num\{0.26\} \&
\textbackslash num\{15.6\} \textbackslash\textbackslash{}
\textbackslash end\{tabular\} \textbackslash end\{ruledtabular\}
\textbackslash end\{table\}

\textbackslash begin\{table\}{[}ht{]} \textbackslash caption\{Mass,
height, and potential energy for Sad Poppers.\}
\textbackslash label\{tab:sad\} \textbackslash begin\{ruledtabular\}
\textbackslash begin\{tabular\}\{cccc\} Trial \& Mass
(\textbackslash unit\{\textbackslash kilo\textbackslash gram\}) \&
Height (\textbackslash unit\{\textbackslash meter\}) \& GPE
(\textbackslash unit\{\textbackslash milli\textbackslash joule\})
\textbackslash\textbackslash{} \textbackslash colrule 1 \& 0.0050 \&
1.63 \& 79.9 \textbackslash\textbackslash{} 2 \& 0.0055 \& 1.30 \& 70.1
\textbackslash\textbackslash{} 3 \& 0.0056 \& 1.53 \& 84.0
\textbackslash\textbackslash{} 4 \& 0.0057 \& 1.72 \& 96.2
\textbackslash\textbackslash{} 5 \& 0.0055 \& 1.58 \& 85.2
\textbackslash\textbackslash{} \textbackslash colrule Mean \&
\textbackslash num\{0.0055\} \& \textbackslash num\{1.55\} \&
\textbackslash num\{83.1\} \textbackslash\textbackslash{} Std. Dev. \&
\textbackslash num\{0.0002\} \& \textbackslash num\{0.16\} \&
\textbackslash num\{9.2\} \textbackslash\textbackslash{}
\textbackslash end\{tabular\} \textbackslash end\{ruledtabular\}
\textbackslash end\{table\}

The results of the two-sample \$t\$-tests are presented in
\textbackslash cref\{tab:stats\}. At the \$\textbackslash alpha = 0.05\$
significance level, no statistically significant differences were found
between the two groups for any of the measured variables.

\textbackslash begin\{table\}{[}ht{]} \textbackslash caption\{Two-sample
\$t\$-test results for Mass, Height, and GPE.\}
\textbackslash label\{tab:stats\} \textbackslash begin\{ruledtabular\}
\textbackslash begin\{tabular\}\{lcccc\} Variable \& \$t\$-value \&
\$df\$ \& \$p\$-value \& \$\textbackslash alpha\$
\textbackslash\textbackslash{} \textbackslash colrule Mass \& 1.74 \& 8
\& 0.12 \& 0.05 \textbackslash\textbackslash{} Height \& 1.27 \& 8 \&
0.24 \& 0.05 \textbackslash\textbackslash{} GPE \& 1.62 \& 8 \& 0.15 \&
0.05 \textbackslash\textbackslash{} \textbackslash end\{tabular\}
\textbackslash end\{ruledtabular\} \textbackslash end\{table\}

\textbackslash begin\{figure\}{[}ht{]} \textbackslash begin\{center\}
\textbackslash includegraphics{[}width=0.45\textbackslash textwidth{]}\{figures/myplot.png\}
\textbackslash end\{center\} \textbackslash caption\{Mean maximum jump
height for happy and sad poppers. Error bars represent
\$\textbackslash pm 1\$ standard deviation (\$n = 5\$ per group).\}
\textbackslash label\{fig:height\_graph\} \textbackslash end\{figure\}
\textbackslash section\{Discussion\} The results of this experiment
indicate that popper design does not have a statistically significant
effect on maximum jump height. Although the mean height of the happy
poppers (\textbackslash qty\{1.72\}\{\textbackslash meter\}) was
slightly greater than that of the sad poppers
(\textbackslash qty\{1.55\}\{\textbackslash meter\}), the distributions
overlap substantially, as indicated by
\textbackslash cref\{fig:height\_graph\}. The two-tailed \$t\$-test gave
a \$p\$-value of \$0.24\$, which is greater than
\$\textbackslash alpha=0.05\$. Therefore, we fail to reject the null
hypothesis, indicating that facial expression does not significantly
affect jump height.

To better analyze the physics of the situation, GPE was calculated for
each trial as well, as seen in \textbackslash cref\{tab:happy\} and
\textbackslash cref\{tab:sad\}. Through this calculation, we were able
to observe the transfer of elastic energy, from the popper being
compressed, to motion (kinetic energy), to height.

The calculated mean GPE values, seen in \textbackslash cref\{tab:happy\}
and \textbackslash cref\{tab:sad\}, show a slight difference of
\textbackslash qty\{13.1\}\{\textbackslash milli\textbackslash joule\}.
However, as seen in \textbackslash cref\{tab:stats\}, a \$t\$-test
between the two energy values resulted in \$p=0.15 \textgreater{}
0.05\$, indicating that the difference is not statistically significant
and likely due to random variation. Since mass influences the conversion
of energy into vertical motion, seen by
\textbackslash cref\{eq:epe,eq:ke\}, a statistical comparison of popper
masses was performed. The comparison showed no statistically significant
difference, with \$p=0.12 \textgreater{} 0.05\$, indicating that mass
was likely not a confounding variable in the height comparison.

The motion of the poppers is consistent with qualitative expectations
from energy conservation. As illustrated by the experimental setup in
\textbackslash cref\{fig:setup\}, compressing the popper requires work
to be done on the material, storing elastic potential energy. Upon
release, this energy is converted into kinetic energy and then into
gravitational potential energy at the peak of the jump. Although
mechanical energy is not perfectly conserved due to non-conservative
forces such as air resistance, sound, and internal deformation of the
rubber, both popper types experienced similar conditions throughout the
experiment. The comparable height distributions observed in
\textbackslash cref\{fig:height\_graph\} therefore suggest that
differences in facial design do not meaningfully affect the efficiency
of energy transfer.

Overall, the statistical results and graphical evidence collectively
support the hypothesis that facial expression does not influence maximum
jump height. Instead, the motion of the poppers is governed primarily by
elastic properties, mass and energy conservation and transfer
mechanisms, rather than surface-level design features.

\textbackslash subsection\{Source of Experimental Error\} A possible
source of experimental error is the force at which each popper is
launched. Despite the fixed location, since each jumper was popped by
one individual and not a repeatable machine, the angle, speed, pressure,
and force that each jumper is pressed on can differ each time, causing
variability in the data. Variability could also have been caused by
error in human reaction. From the videos of each popper, the max height
was analyzed. Humans can misread exact measurements on a meter stick,
which could cause inconsistencies in data. Finally, the most likely
cause of error is wear. As the poppers were used before, it is likely
that there were internal deformations that may have led to
inconsistencies within the data.

\textbackslash section\{Acknowledgements\} K. Coulanges, S. Dakwale, and
S. Hein primarily led on data collection and recording. S. Hein
compressed all the poppers to be released and wrote the abstract. S.
Yellapragada wrote the introduction, methods and materials, R code for
the graph, and revised the entire lab manuscript. R. Wallace calculated
results based on height data. S. Dakwale explained sources of
experimental error and performed all statistical analyses. K. Coulanges
created setup figures and wrote the discussion. All members contributed
to proofreading and editing, as well as feedback on the overall paper.
We would like to thank Dr. Evangelista for providing guidance,
materials, and supervision throughout the lab.

\textbackslash bibliography\{example.bib\}
\textbackslash end\{document\}
